\section{Introduction}
\label{sec:introduction}

For a partition $\alpha\vdash n$, let $S^\alpha$ denote the corresponding
irreducible complex representation of the symmetric group $S_n$.
The Kronecker coefficient
\[
g(\alpha,\beta,\lambda)
=\dim\Hom_{S_n}(S^\lambda,S^\alpha\otimes S^\beta)
\]
is the multiplicity of $S^\lambda$ in the tensor product, with $S_n$
acting diagonally. Put
\begin{equation}\label{intro:staircase}
N_m=\frac{m(m+1)}2,\qquad \rho_m=(m,m-1,\ldots,1).
\end{equation}
We prove the following statement.

\begin{theorem}[Saxl's conjecture]\label{thm:saxl}
For every integer $m\ge1$ and every partition $\lambda\vdash N_m$,
\[
g(\rho_m,\rho_m,\lambda)>0.
\]
Equivalently, $S^{\rho_m}\otimes S^{\rho_m}$ contains every irreducible
complex representation of $S_{N_m}$.
\end{theorem}

The assertion is uniform in both the staircase size and the constituent.
The proof establishes it inside one explicitly generated submodule of
the tensor square, a stronger conclusion that retains the vector needed
at each induction step.

The companion \cite[Theorem~1.1]{universaltensorsquares2026} uses the
stronger cyclic statement of \Cref{thm:cyclic-saxl}, together with
additional arguments, to construct, for every positive integer
$n\notin\{2,4,9\}$, an irreducible complex representation of $S_n$ whose
tensor square contains every irreducible representation of $S_n$.

\subsection{History and significance}

Saxl's conjecture belongs to a broader tensor-square covering problem:
can a single irreducible representation have every irreducible as a
constituent of its square? A motivating precedent comes from finite
simple groups of Lie type. Heide, Saxl, Tiep, and Zalesski proved that
the Steinberg square has this property, except for the groups
$\mathrm{PSU}_d(q)$ with $d\ge3$ and $\gcd(d,2(q+1))=1$
\cite[Theorem~1.2]{heidestz2013}.
Saxl proposed the staircase assertion in a UCLA Combinatorics Seminar on
20 March 2012, as recorded by Pak, Panova, and Vallejo
\cite[Conjecture~1.2 and Footnote~2]{pakpanovavallejo2016}.
Their work places the staircase square among tensor-square covering
questions for symmetric groups and proves several families of positive
Kronecker coefficients; in particular, their Theorem~1.4 covers hooks
and two-row partitions for sufficiently large staircase size.
Ikenmeyer proved positivity when $\lambda$ and $\rho_m$ are comparable
in the dominance order \cite[Theorem~2.1]{ikenmeyer2015}.
We use his triangular arrangement of two alternating tensor layers
\cite[Section~4]{ikenmeyer2015}; his argument also yields all hook constituents
\cite[Corollary~6.1]{ikenmeyer2015}.

Character-theoretic and modular methods have supplied further substantial
families. The character-value criterion of Pak, Panova, and Vallejo
\cite{pakpanovavallejo2016} detects constituents using the principal-hook
cycle type of a self-conjugate partition. Bessenrodt developed related
critical-class and spin-character methods, including the double-hook
case \cite{bessenrodt2018}. Li proved the triple-hook case using
computer-assisted finite reductions \cite{li2021}.
Here double and triple hooks are diagrams whose largest square has side
two and three, respectively.
Bessenrodt, Bowman, and Sutton connected Kronecker positivity to
2-modular decomposition numbers; their results include all
2-height-zero constituents, hence all odd-degree constituents, and
framed staircases \cite[Theorems~5.2 and~5.5]{bessenrodtbowmansutton2021}.

Another line of work addresses the tensor exponent and asymptotic
coverage. Luo and Sellke used the semigroup property to prove that a
random partition occurs in the staircase square with probability tending
to one, for both the uniform and Plancherel measures
\cite[Theorems~1.5 and~1.6]{luosellke2017}. They also found, in every
sufficiently large degree, an irreducible representation whose fourth
tensor power contains every irreducible
\cite[Theorem~1.4]{luosellke2017}. Harman and Ryba proved full coverage
for the tensor cube of every staircase, with combinatorial and modular
proofs \cite[Theorem~1.1]{harmanryba2023}.
Asymptotic coverage does not exclude exceptional constituents, while a
cube or fourth power uses additional tensor factors; neither gives the square
assertion in \Cref{thm:saxl}.

Recent developments include semigroup constructions of further
constituents and generalized-block constraints on support
\cite{ebrahimi2025}, and a full
staircase-square theorem for unipotent characters of $\GL_n(q)$
\cite[Theorem~1.1]{letelliernam2025}. The latter is an analogue in a
different character theory: ordinary Kronecker positivity implies the
corresponding unipotent positivity, but the converse need not hold.

The band calculation below also has a precise predecessor at the level
of ambient representations. Brown, van Willigenburg, and Zabrocki proved
that
$S^{(m,m-1)}\otimes S^{(m,m-1)}$ is the multiplicity-free sum of all
irreducibles of $S_{2m-1}$ indexed by partitions with at most four rows
\cite[Corollary~4.1]{brownwilligenburgzabrocki2010}.
Our induction needs each of these constituents in the cyclic module of
one prescribed band vector, rather than merely somewhere in that
ambient square. Explicit nonzero pairings supply this additional
generator-specific statement. Throughout, the objective is support,
or positivity, rather than a formula for Kronecker multiplicities.

\subsection{The construction and its reusable steps}

Place $N_m$ tensor positions in the triangle
$B_m=\{(i,j):i,j\ge1,\ i+j\le m+1\}$.
Let $R_m$ and $C_m$ be its row and column set partitions, ordered within
each row by increasing $j$ and within each column by increasing $i$.
Use the ordered basis $e_1,\ldots,e_m$ of $\C^m$. On a block of length
$r$, alternate the initial letters $e_1,\ldots,e_r$, taking the signed
sum over all permutations. Taking the product over row blocks gives
$v_{R_m}$; doing the same over column blocks gives $v_{C_m}$.
These are polytabloids of staircase type. Instead of beginning with an
arbitrary vector in the tensor square, we work throughout with
\[
W_m=\C[S_{B_m}](v_{R_m}\otimes v_{C_m}),
\]
where $S_{B_m}$ permutes the positions in both layers simultaneously.
\Cref{thm:cyclic-saxl} asserts that this particular cyclic module already
contains every irreducible. We first prove that it contains every
$S^\lambda$ for which $\lambda$ is dominated by $\rho_m$: each initial
sum of the parts of $\lambda$ is at most the corresponding staircase
sum. A sign twist and a permutation-module quotient give this starting
case (\Cref{prop:dominated}). The following three steps reach the
remaining partitions.

\begin{enumerate}
\item \textbf{A full induced quotient from a coordinate projection.}
Separate a smaller triangle $B_{m-s}$ from a band of width $s=1$ or $2$.
Call letters $1,\ldots,s$ low and letters $s+1,\ldots,m$ high.
Each position carries one letter from each tensor layer.
Project onto words whose two letters at each position are either both
high or both low. The staircase row and column counts force the high
positions of the projected generator to be exactly the smaller triangle.
The generator then factors into the smaller cyclic generator and a band
generator, whose cyclic module is denoted by $U_{m,s}$.
Translates of the high-position set occupy disjoint
coordinate sectors, so the quotient is the full induced module
\[
\Ind_{S_{B_{m-s}}\times S_{B_m\setminus B_{m-s}}}^{S_{B_m}}
       (W_{m-s}\boxtimes U_{m,s}).
\]
\Cref{lem:sector-induction} isolates the general reason that sector
separation gives an induced-module isomorphism, including an explicit
inverse. \Cref{prop:band-quotient} gives the resulting quotient.

\item \textbf{Constituents detected in the specified band generator.}
A width-two band is an odd path of alternating pairs with fixed endpoint
letters. At each position, a covector on the pair-letter space is
represented by a two-by-two matrix. Pairing along the path evaluates a word of
matrices and adjugates. To test a partition with at most four rows,
place its tableau columns on consecutive path segments and alternate
the covectors within each column. Each segment has length at most four.
A four-element basis of $\Mat_2(\C)$ makes every such alternated segment
matrix invertible.
A simultaneous conjugation of the test basis makes the endpoint entry
of the product nonzero. The resulting dual-polytabloid pairing proves
that every shape with at most four rows occurs in $U_{m,2}$ itself
(\Cref{prop:two-band}).

\item \textbf{A horizontal-strip reduction with four steps.}
If a shape outside this dominance range has a first row of length at
least $m$, one horizontal $m$-strip suffices. Otherwise its first four
rows contain more than $2m-1$ boxes. Successive sweeps of the bottom
boxes of columns then remove exactly $2m-1$ boxes in at most four
horizontal strips (\Cref{comb:strip-reduction}). In this second case,
let $\nu$ be the remaining partition. Iterated Pieri supplies a shape
$\eta\vdash2m-1$ with at most four rows such that $S^\lambda$ occurs
in the representation induced from $S^\nu\boxtimes S^\eta$.
The band calculation puts $S^\eta$ in $U_{m,2}$, while induction puts
$S^\nu$ in $W_{m-2}$. The full induced quotient then gives the target
constituent in $W_m$.
\end{enumerate}

The smaller factor in the quotient is the actual cyclic module $W_{m-s}$,
and the band factor is generated by the actual projected vector.
This is why both parts of the induction track generators instead of
using ambient tensor-product support alone. Only one copy of a suitable
intermediate irreducible is needed.

\Cref{sec:specht-tools} fixes the tensor conventions and proves the
needed form of Young's rule from Pieri. \Cref{sec:dominance} defines
$W_m$ and proves the dominated case. The quotient and band calculations
occupy \Cref{sec:band-quotient,sec:path-band};
\Cref{sec:completion} proves the strip reduction and both main theorems.
